% ============================================================
% 第 3 章 变压器概述（教材第二篇，印刷页 24--29）
% 节名：第一节 变压器的分类；第二节 变压器的基本结构；第三节 变压器的型号与额定值
% ============================================================
\chapter{变压器概述}

变压器是电力系统输电与配电的主要设备，也是测量、控制及通信等装置的重要元件，被广泛应用在各个领域，因此种类很多。本章只解决三个问题：有哪些变压器（分类）、它由什么构成（结构）、它的容量电压电流是怎么规定的（额定值）。这三件事是后面所有分析计算的"边界条件"——等值电路里的 $U_{1N}$、$I_{1N}$、$S_N$ 就是从这里来的。

\section{变压器的分类}

\subsection{按用途分类}
\begin{enumerate}
  \item \textbf{电力变压器}：在电力系统中传送和分配电能。容量大（从几十千伏安到百万伏安），电压高（从几百伏到兆伏），应用范围广。
  \item \textbf{调压变压器}：在电网或实验室中用于调节电压。
  \item \textbf{测量变压器}：如电流互感器和电压互感器等。
  \item \textbf{专用变压器}：如电炉变压器、电焊变压器、整流变压器、高压试验变压器，以及无线电通信、自动控制系统和各种仪器用的小功率变压器等。
\end{enumerate}

\subsection{按其他方式分类}
\begin{itemize}
  \item 按\textbf{绕组数目}：双绕组变压器、三绕组变压器等；
  \item 按\textbf{铁芯结构}：心式变压器、壳式变压器；
  \item 按\textbf{相数}：单相变压器、三相变压器和多相变压器（如电力电子装置中的六相、十二相变压器）；
  \item 按\textbf{冷却方式}：油浸式变压器、干式变压器。
\end{itemize}

变压器的种类虽多，但基本工作原理一致。本书着重讨论\textbf{电力变压器}。

\begin{center}
\begin{tabular}{@{}llp{7.6cm}@{}}
\toprule
分类依据 & 类别 & 特点与典型用途 \\
\midrule
用途 & 电力变压器 & 输电与配电，容量几十 kV·A～百万 kV·A \\
 & 调压变压器 & 电网或实验室调压 \\
 & 测量变压器 & 电流互感器、电压互感器 \\
 & 专用变压器 & 电炉、电焊、整流、高压试验、小功率通信与控制 \\
铁芯结构 & 心式 & 绕组包围铁芯柱，结构简单、绝缘易做，应用最广 \\
 & 壳式 & 铁芯包围绕组，机械强度高但制造复杂，用于低压大电流 \\
绕组数目 & 双/三绕组 & 三绕组用于三种电压等级的交汇点 \\
相数 & 单相/三相/多相 & 三相最常用；多相用于电力电子装置 \\
冷却 & 油浸式/干式 & 油浸式散热好、容量大；干式防火性能好 \\
\bottomrule
\end{tabular}
\end{center}

\section{变压器的基本结构}

变压器通过套装在同一铁芯柱上的两个或两个以上的绕组工作，因此基本结构件是\textbf{铁芯}和\textbf{绕组}。对于油浸式电力变压器，还有油箱、变压器油、绝缘套管和其他结构附件。铁芯和绕组合称为变压器的\textbf{器身}。

\subsection{铁芯}
铁芯是变压器的\textbf{磁路}，又是变压器的\textbf{机械骨架}。它由\textbf{铁芯柱}和\textbf{铁轭}两部分组成：铁芯柱上套装绕组，铁轭把各个铁芯柱连接起来，形成闭合铁芯磁路。

为减少铁芯中的磁滞和涡流损耗，铁芯通常选用厚度为 0.23、0.35\,mm（目前多为 0.30\,mm 及以下）的两面涂绝缘漆的\textbf{硅钢片}叠成。硅钢片的铁芯损耗小、磁导率高。

\subsubsection{心式与壳式}
\begin{itemize}
  \item \textbf{心式}：绕组包围铁芯柱。结构简单，绕组套装和绝缘比较容易，应用比较广泛。
  \item \textbf{壳式}：铁芯包围绕组。铁芯机械强度高，但制造复杂，通常只有低压大电流的变压器或小容量的电信变压器采用。
\end{itemize}

\subsubsection{叠片与接缝}
铁芯叠片通常将硅钢片剪成一定尺寸的长方形，再按要求叠装成铁芯。叠片\textbf{交错叠装}，使相邻两层的接缝互相错开，这样叠成的铁芯接缝气隙小、导磁性能好。

硅钢片有热轧和冷轧两种：
\begin{itemize}
  \item \textbf{热轧硅钢片}导磁方向性不明显，叠积铁芯可采用\textbf{直接缝}；
  \item \textbf{冷轧硅钢片}导磁性更好、损耗更小，变压器几乎都采用冷轧硅钢片；但它具有很强的\textbf{方向性}，为避免磁通转向时在拐角处增加磁阻，通常采用\textbf{斜接缝}。
\end{itemize}
由于直接缝比斜接缝的材料利用率高，变压器中直接缝与斜接缝混合使用也较普遍。

\subsubsection{铁芯柱截面}
小型变压器铁芯柱的横截面是\textbf{正方形}；容量较大的变压器，为充分利用绕组的内圆空间，一般采用\textbf{多级阶梯形}截面。当铁芯柱横截面外接圆直径大于 500\,mm 时，中间还应留出\textbf{油道}以改善铁芯内部的散热条件。铁轭的横截面形状一般为矩形，面积比铁芯柱的横截面积大 5\%～10\%，以减少励磁电流和铁芯损耗。

\subsection{绕组}
绕组是变压器的\textbf{电路部分}，通常用纸包的绝缘扁导线或漆包的圆导线绕制而成，导线多为铜线。根据高、低压绕组在铁芯柱上的相对位置，绕组的排列方式分为两种。

\subsubsection{同心式}
高、低压绕组同心地套装在铁芯柱上，\textbf{低压绕组在里面，高压绕组在外面}。这样可使绝缘距离小，绕组和铁芯的尺寸都可以做得小些，同时也便于引出高压绕组的分接头。同心式绕组制造方便，国产变压器多采用这种结构。

\subsubsection{交叠式}
高、低压线圈都做成饼状，沿铁芯柱高度方向交替放置；靠近上下铁轭的位置一般放置低压绕组。交叠式绕组\textbf{漏抗较小，机械强度较高}，易于接成多条并联支路，主要用于低电压、大电流的电焊变压器和电炉变压器。

绕组有电流流通时要发热，中大型变压器的绕组之间、一个绕组的各层之间，有时甚至各匝之间都垫上绝缘块，形成\textbf{油路}以利散热。

在配电变压器中，高压绕组一般都引出几个\textbf{分接头}，在电网电压有波动时，可通过分接开关接通不同的分接头，调节高压绕组的匝数，改变变压器的实际匝数比，使低压侧输出电压保持相对稳定。

变压器输入电能一侧的绕组称为\textbf{一次绕组}，输出电能一侧的绕组称为\textbf{二次绕组}。两侧匝数不等，通过电磁感应作用，一次绕组把电能传递到二次绕组，二次绕组就能得到不同的电压和电流。

\begin{center}
\begin{tabular}{@{}p{1.7cm}p{4.4cm}p{3.3cm}p{3.6cm}@{}}
\toprule
排列方式 & 结构 & 优点 & 主要用途 \\
\midrule
同心式 & 低压在内、高压在外，同心套在铁芯柱上 & 绝缘距离小、尺寸小、便于引出分接头、制造方便 & 国产电力变压器普遍采用 \\
交叠式 & 高、低压饼状线圈沿柱高交替放置，靠铁轭处放低压 & 漏抗小、机械强度高、易多支路并联 & 低电压大电流（电焊、电炉变压器） \\
\bottomrule
\end{tabular}
\end{center}

\subsection{油箱及变压器油}
油箱是变压器的外壳，由钢板焊接而成。油箱内盛变压器油，器身浸在油中。变压器油有\textbf{绝缘}和\textbf{散热}两个作用：绕组和铁芯产生的热量由变压器油传递给油箱壁、散热管或散热器，从而冷却器身。

\subsection{绝缘套管}
绝缘套管是变压器绕组的引出线之间以及引出线与油箱之间的绝缘，并起固定引出线的作用，用陶瓷制成，有多种形式。

除此之外，油浸式电力变压器中通常还有\textbf{储油柜}、\textbf{散热器}、\textbf{气体继电器}、\textbf{安全气道}、\textbf{吸湿器}、\textbf{分接开关}、\textbf{油位计}和继电保护装置等附件。

\subsection{干式变压器}
自 20 世纪 80 年代开始，干式变压器在我国得到迅速发展。干式变压器的铁芯和绕组不浸在任何液体中，具有运行安全可靠、维护简单、可深入负荷中心等优点。但由于空气的绝缘强度和散热性能都比油差，以空气绝缘的干式变压器有效材料消耗比油浸式多，承受冲击电压的能力也比油浸式差。干式变压器主要适用于高层建筑、机场和地铁等安全防火要求较高的配电场所，主要由铁芯和绕组组成，铁芯及其夹件充当机壳；绕组通常用纱包的绝缘扁铜导线绕制，为增强绝缘性能通常进行\textbf{树脂浇注}；较大容量的干式变压器有时配装风冷装置和罩壳。

\corepoint{结构与原理的两条对应关系必须记牢：① 铁芯为什么要闭合、要叠片、要斜接缝——都是为了"用最小的励磁磁动势建立主磁通，同时把铁耗压到最低"；② 低压绕组为什么放里面——为了减小高压绕组对铁芯的绝缘距离，从而缩小整体尺寸。凡是"结构为什么这样"的问题，答案通常落在\textbf{绝缘、散热、损耗、机械强度}这四个词上。}

\section{变压器的型号与额定值}

\subsection{变压器的型号}
变压器的型号包括说明结构特点的基本代号、额定容量和额定电压。例如型号 \textbf{S-500/10} 表示：三相油浸自冷式、双绕组电力变压器，额定容量 500\,kV·A，额定电压 10\,kV。

\begin{center}
\begin{tabular}{@{}ll@{}}
\toprule
代号 & 含义 \\
\midrule
S & 三相（D 为单相） \\
油浸自冷 & 冷却方式（干式、强迫油循环风冷等另有代号） \\
双绕组 & 绕组数目 \\
500 & 额定容量，kV·A \\
10 & 高压侧额定电压等级，kV \\
\bottomrule
\end{tabular}
\end{center}

\subsection{额定值}
变压器的额定值是制造厂对变压器正常运行时所作的使用规定，也是设计、使用和试验变压器的依据。在额定状态下运行，可以保证变压器长期可靠地工作，并具有优良的性能。

\subsubsection{额定容量 $S_N$}
在规定的额定状态下变压器输出的\textbf{视在功率}，单位 kV·A。对于双绕组电力变压器，一、二次绕组的额定容量设计值相同。

\subsubsection{额定电压 $U_{1N}$ 和 $U_{2N}$}
在规定的额定状态下，一次绕组所加的电压值称为一次绕组的额定电压 $U_{1N}$；当一次绕组加额定电压 $U_{1N}$ 时，二次绕组所得到的\textbf{空载电压}称为二次绕组的额定电压 $U_{2N}$。单位 V 或 kV。\textbf{三相变压器的额定电压是指线电压}。

\subsubsection{额定电流 $I_{1N}$ 和 $I_{2N}$}
根据额定容量和额定电压算出的\textbf{线电流}，单位 A。

\keyformula{\text{单相：}I_{1N}=\frac{S_N}{U_{1N}},\quad I_{2N}=\frac{S_N}{U_{2N}};\qquad
\text{三相：}I_{1N}=\frac{S_N}{\sqrt3\,U_{1N}},\quad I_{2N}=\frac{S_N}{\sqrt3\,U_{2N}}}

\subsubsection{额定频率 $f_N$}
我国额定工频为 50\,Hz。此外，额定工作状态下变压器的效率、温升等数据均属额定值。由于额定值都标注在变压器外壳的铭牌上，所以额定值也称为\textbf{铭牌值}。

\begin{center}
\begin{tabular}{@{}lll@{}}
\toprule
额定值 & 定义要点 & 单位 \\
\midrule
$S_N$ & 额定状态下输出的视在功率；双绕组一、二次侧设计值相同 & kV·A \\
$U_{1N}$ & 一次绕组应加的电压（三相指线电压） & V 或 kV \\
$U_{2N}$ & 一次加 $U_{1N}$ 时二次绕组的\textbf{空载}电压（三相指线电压） & V 或 kV \\
$I_{1N}$、$I_{2N}$ & 由 $S_N$ 与 $U_N$ 算出的线电流 & A \\
$f_N$ & 我国工频 50\,Hz & Hz \\
\bottomrule
\end{tabular}
\end{center}

\corepoint{三个易错点：① $U_{2N}$ 是\textbf{空载}电压，不是额定负载下的端电压——带负载后二次电压会因漏阻抗压降而低于 $U_{2N}$，这正是第 6 章"电压变化率"要算的东西；② 三相变压器的 $U_N$、$I_N$ 都是\textbf{线值}，换算到相值要除以 $\sqrt3$（星接）或乘以 $\sqrt3$（角接），做参数计算时务必先统一到相值；③ $S_N$ 是视在功率，不是有功功率，故 $I_N=S_N/(\sqrt3U_N)$ 与负载功率因数无关。}
